\subsection{nml\_grid}
\label{SEC:nml_grid}
\inp*[nmlGrid]{}[nml\_grid][]

This \hyperref[SEC:User_Input]{namelist block}
is used to specify parameters related to the grid input file.

\begin{nmlblocklst}[
,label={LST:nml_grid}
,caption={Namelist nml\_grid.}
]
&nml_grid
  Grid_Scaling_Factor      =   1.0
  Grid_Format              =  -1
  Grid_File                =  ""
  Plot3d_Cutfile           =  ""
  Plot3d_BCSfile           =  ""
  Plot3d_Stride_I          =  -1
  Plot3d_Stride_J          =  -1
  Plot3d_Stride_K          =  -1
  Plot3d_Stride_All        =  -1
  Plot3d_Agglomerate_Order =  -1
  Correct_CGNS_Grid_Points =   F
  Internal_Grid_DoF        =   0
  Internal_Grid_Order      =  -1
  Metis_Option             =   3
  Geom_Wts(1)%Tria         =   1
  Geom_Wts(1)%Quad         =   1
  Geom_Wts(1)%Tetr         =   1
  Geom_Wts(1)%Pyra         =   1
  Geom_Wts(1)%Pris         =   1
  Geom_Wts(1)%Hexa         =   1
/
\end{nmlblocklst}
%%Geom_Wts(1)%Edge         =   1
%%Geom_Wts(1)%Face_Tria    =   1
%%Geom_Wts(1)%Face_Quad    =   1

%==========================================================================

\begin{nmldescitem}
Grid_Scaling_Factor = 1.0
\end{nmldescitem}
\inp*[nmlGrid::GridScalingFactor]{}[Grid\_Scaling\_Factor][]

\begin{nmldescbody}
Multiplication factor ($> 0.0$) to scale the input grid file coordinates to
the desired length units in meters.
%
For example, use a value of 0.0254 to convert a grid containing coordinates
in inches to meters.
%
See the discussion of
  \key[key:GridUnits]{grid units}[grid units]
and
  \key[key:GridScaledUnits]{grid scaled units}[grid scaled units]
in \sref{SEC:Grid_Scaling}.
\end{nmldescbody}

%==========================================================================

\begin{nmldescitem}
Grid_Format = -1
\end{nmldescitem}
\inp*[nmlGrid::GridFormat]{}[Grid\_Format][]

\begin{nmldescbody}
Specifies the format of the input grid file.

\begin{tabular}{ @{} r @{ = } p{0.85\linewidth} @{} }
  -1 & Internally generated grid
       (only supported with certain 
       \inp[nmlInitialization::Initialization]{Initialization} options) \\
   1 & Gambit Neutral (very limited support, high-order curved elements not
       supported) \\
   2 & Gmsh (limited support) \\
   3 & Plot3D \\
   4 & CGNS (recommended)
\end{tabular}

\end{nmldescbody}

%==========================================================================

%##%\FINALFIX{\newpage}
\begin{nmldescitem}
Grid_File = ""
\end{nmldescitem}
\inp*[nmlGrid::GridFile]{}[Grid\_File][]

\begin{nmldescbody}
File path to the input grid file relative to the current directory.
\end{nmldescbody}

%==========================================================================

\begin{nmldescitem}
Plot3d_Cutfile = ""
\end{nmldescitem}
\inp*[nmlGrid::Plot3dCutfile]{}[Plot3d\_Cutfile][]

\begin{nmldescbody}
File path to the input file that specifies the cut conditions for a
multiblock Plot3D grid file.
\end{nmldescbody}

%==========================================================================

\begin{nmldescitem}
Plot3d_BCSfile = ""
\end{nmldescitem}
\inp*[nmlGrid::Plot3dBCSfile]{}[Plot3d\_BCSfile][]

\begin{nmldescbody}
File path to the input file that specifies the boundary conditions for a
Plot3D grid file.
\end{nmldescbody}

%==========================================================================

\begin{nmldescitem}
Plot3d_Stride_All = -1
Plot3d_Stride_I   = -1 
Plot3d_Stride_J   = -1
Plot3d_Stride_K   = -1
\end{nmldescitem}
\inp*[nmlGrid::Plot3dStrideAll]{}[Plot3d\_Stride\_All][]%
\inp*[nmlGrid::Plot3dStrideI]{}[Plot3d\_Stride\_I][]%
\inp*[nmlGrid::Plot3dStrideJ]{}[Plot3d\_Stride\_J][]%
\inp*[nmlGrid::Plot3dStrideK]{}[Plot3d\_Stride\_K][]%
 
\begin{nmldescbody}
If $> 0$, applies a coarsening stride to the specified direction
of all Plot3D blocks.
 
\inptext{Plot3d\_Stride\_All} is applied first to all directions.
%
Then the individual i, j, k options can be used to override the
stride setting in that specific direction.
\end{nmldescbody}

%==========================================================================

\begin{nmldescitem}
Plot3d_Agglomerate_Order = -1
\end{nmldescitem}
\inp*[nmlGrid::Plot3dAgglomerateOrder]{}[Plot3d\_Agglomerate\_Order][]

\begin{nmldescbody}
If $> 0$, agglomerates linear Plot3D cells into high-order grid cells of
this polynomial order.
%
For this to be successfull, the number of {\em cells} in each direction
for every block must be divisible by this value.
\end{nmldescbody}

%==========================================================================

\begin{nmldescitem}
Correct_CGNS_Grid_Points = F
\end{nmldescitem}
\inp*[nmlGrid::CorrectCGNSGridPoints]{}[Correct\_CGNS\_Grid\_Points][]

\begin{nmldescbody}
Experimental code that tries to correct the high-order points given in a
CGNS grid file created by Pointwise. The high-order points should be equally
spaced, but some software such as Pointwise does not require this.
\end{nmldescbody}

%==========================================================================

\begin{nmldescitem}
Internal_Grid_DoF = 0
\end{nmldescitem}
\inp*[nmlGrid::InternalGridDoF]{}[Internal\_Grid\_DoF][]
\key*{}[degree-of-freedom (DoF)!internally generated grid]
\key*{}[internally generated grid!degree-of-freedom (DoF)]

\begin{nmldescbody}
This specifies the grid resolution to use for an
\key[key:InternallyGeneratedGrid]{internally generated grid}[=].
\\
\begin{tabular}{ @{\quad} r @{ \quad } p{0.85\linewidth} @{} }
  $<0$ & User specified number of cells. \\
  $=0$ & 512 degrees of freedom. \\
  $>0$ & User specified degrees of freedom.
\end{tabular}\\
%
If the number of degrees of freedom are specified, then the number of cells
needed to achieve the specified degrees of freedom in each direction will be
computed from the solution order as:
\begin{equation*}
    \left[ \text{\# of cells} \right]
  = \mathrm{nint}
    ( \inptext{Internal\_Grid\_DoF}
%   ( \text{\inp[nmlGrid::InternalGridDoF]{Internal\_Grid\_DoF}}
    / ( \text{\inp[nmlNumerics::SolutionOrder]{Solution\_Order}} + 1 )
    ).
\end{equation*}

\NOTE
An internally generated grid will only be created for specific
\inp[nmlInitialization::Initialization]{Initialization} conditions,
as indicated on \pref{\pageref{nmlInitialization::Initialization}}.
%
For the following cases, one should instead use the grid inputs in the
associated namelist block.

\begin{tabular}{S ll}
  {\inp[nmlInitialization::Initialization]{Initialization}}
& Description
& Namelist Block \\
  \hline
   2 &  Channel flow
     &  \inp[nmlTestcaseChannel]{nml\_testcase\_channel} \\
  21 &  Taylor-Couette flow
     &  \inp[nmlTestcaseTCF]{nml\_testcase\_tcf} \\
  22 &  Freestream Taylor-Couette flow
     &  \inp[nmlTestcaseTCF]{nml\_testcase\_tcf}
\end{tabular}
\end{nmldescbody}

%==========================================================================

\begin{nmldescitem}
Internal_Grid_Order = -1
\end{nmldescitem}
\inp*[nmlGrid::InternalGridOrder]{}[Internal\_Grid\_Order][]

\begin{nmldescbody}
Polynomial order used by the internal grid generator to create a
high-order grid.
%
This only applies to internally generated grids with curved geometries
(\thatis Taylor-Couette flow, channel, and annular pipe channel test cases).
%
All other internally generated grids have linear geometries.
\end{nmldescbody}

%==========================================================================

\FINALFIX{\newpage}
\begin{nmldescitem}
Metis_Option = 3
\end{nmldescitem}
\inp*[nmlGrid::MetisOption]{}[Metis\_Option][]

\begin{nmldescbody}
Specifies the \key{METIS}[=] algorithm to use for
partitioning the global grid. If the specified option fails, it will
loop through all of the other options to try to find a successful partitioning.

%% Note: longtable permits page breaks.
%\begin{tabular}{ @{} r @{ = } p{0.85\linewidth} @{} }
\begin{longtable}{ @{} r @{ = } p{0.85\linewidth} @{} }
  1 & Create non-weighted partitions based on a multilevel recursive bisection
      while minimizing the edge cut.
      \\
  2 & Create non-weighted partitions based on a multilevel $k$-way algorithm
      while minimizing the edge cut.
      \\
  3 & Create non-weighted partitions based on a multilevel $k$-way algorithm
      while minimizing the total communication volume.
      \\
  4 & Create weighted partitions based on a multilevel $k$-way algorithm
      while minimizing the total communication volume.
      This is equivalent to option 3 but also
      includes cell weights (based on geometry type).
      \\
  5 & Create weighted partitions based on a multilevel $k$-way algorithm
      while minimizing the edge cut.
      This is equivalent to option 2 but also
      includes cell weights (based on geometry type).
      \\
  6 & Create weighted partitions based on a multilevel recursive bisection
      while minimizing the edge cut.
      This is equivalent to option 1 but also
      includes cell weights (based on geometry type).
\end{longtable}
%\end{tabular}

Note that the last three options include geometry-dependent
computational cost into load balancing the partitions.
%
Weighting for each geometry type and order is specified through the input
\inp[nmlGrid::GeomWts]{Geom\_Wts}[=][].

\end{nmldescbody}

%==========================================================================

\begin{nmldescitem}
Geom_Wts(1)%Tria = 1
Geom_Wts(1)%Quad = 1
Geom_Wts(1)%Tetr = 1
Geom_Wts(1)%Pyra = 1
Geom_Wts(1)%Pris = 1
Geom_Wts(1)%Hexa = 1
\end{nmldescitem}
%%If including these, then align = signs above.
%Geom_Wts(1)%Edge      = 1
%Geom_Wts(1)%Face_Tria = 1
%Geom_Wts(1)%Face_Quad = 1
\inp*[nmlGrid::GeomWts]{}[Geom\_Wts][]
%\inp*[nmlGrid::GeomWts/Edge]{}[Geom\_Wts(:)\%Edge][]
\inp*[nmlGrid::GeomWts/Tria]{}[Geom\_Wts(:)\%Tria][]
\inp*[nmlGrid::GeomWts/Quad]{}[Geom\_Wts(:)\%Quad][]
\inp*[nmlGrid::GeomWts/Tetr]{}[Geom\_Wts(:)\%Tetr][]
\inp*[nmlGrid::GeomWts/Pyra]{}[Geom\_Wts(:)\%Pyra][]
\inp*[nmlGrid::GeomWts/Pris]{}[Geom\_Wts(:)\%Pris][]
\inp*[nmlGrid::GeomWts/Hexa]{}[Geom\_Wts(:)\%Hexa][]
%\inp*[nmlGrid::GeomWts/FaceTria]{}[Geom\_Wts(:)\%Face\_Tria][]
%\inp*[nmlGrid::GeomWts/FaceQuad]{}[Geom\_Wts(:)\%Face\_Quad][]

\begin{nmldescbody}
These are the geometry weights used by the \key{METIS}[=]
partitioning algorithm to split the mesh between processors.
%
Valid values are integers in the range [0:100], with larger values
indicating those geometries must perform more computational work
relative to geometries with lower values.
%
Note that the \inpstyle{Tria} and \inpstyle{Quad} inputs are for
element volumes for 2D simulations, and thus are not used for 3D simulations.
%
%Note that the \inpstyle{Tria} and \inpstyle{Quad} inputs are for
%element volumes in two dimensions, whereas the \inpstyle{Face\_Tria}
%and \inpstyle{Face\_Quad} inputs are for element surfaces in
%three dimensions.
%%
%Currently, the \inpstyle{Face\_Tria} and \inpstyle{Face\_Quad} weights are not
%used.
%
See the \inp[nmlGrid::MetisOption]{Metis\_Option} input
for additional details.

Preliminary testing of the geometry weights indicates that improvements
in performance may be highly dependent upon the number and distribution
of each cell type throughout the domain and how many partitions are
requested.
%
We therefore offer no specific weighting recommendations.
 
\IDEA{
  Add table of weights for either Rome or Milan nodes since AMD nodes are used
  at both NAS and a lot of DoD machines, just to serve as a rough starting
  point. Some things to probably point out:
  \begin{itemize}
    \item
      Weights are only necessary for unstructured grids.
    \item
      Not using weights can severely degrade performance for large grids using
      large numbers of MPI processes (10k+).
    \item
      Unless the user has compiled a custom METIS library using 64-bit
      integers, the [maximum weight value]*[number of grid cells] needs to be
      smaller than $2^{31}$ or the partitioning will fail. Use one of the
      example Rome\slash Milan weights from the table to illustrate this point
      with an example of too many cells for these weight values, and how to
      scale the weights to get them within the 32-bit integer limit.
    \item
      Possibly provide an example of how unpredictable the weights can be to
      make the user aware of this issue.
    \item
      Include some remarks about using weights with the different partitioning
      methods and how unpredictable the results can be. Maybe include an example
      of this with the example of how unpredictable the weights can be.
  \end{itemize}
}
\end{nmldescbody}

%==========================================================================

